\begin{table}[!ht]\centering
\caption{The occupation-year panel by release, 2014--2020}\label{tab:sample}
\begin{threeparttable}\footnotesize\setlength{\tabcolsep}{3pt}
\begin{tabular}{lcccccc}\toprule
 & & \multicolumn{2}{c}{Mean log hourly pay} & \multicolumn{2}{c}{Cross-section gradient} & Share of jobs in \\
\cmidrule(lr){3-4}\cmidrule(lr){5-6}
Release & Occupations & Below-median & Above-median & Slope & (s.e.) & above-median risk (\%) \\ \midrule
2014 & 364 & 2.891 & 2.360 & -2.114 & (0.083) & 54.0 \\
2015 & 365 & 2.896 & 2.380 & -2.071 & (0.081) & 54.1 \\
2016 & 365 & 2.902 & 2.399 & -2.030 & (0.079) & 52.8 \\
2017 & 363 & 2.924 & 2.419 & -2.009 & (0.080) & 52.6 \\
2018 & 364 & 2.948 & 2.452 & -1.980 & (0.079) & 52.3 \\
2019 & 363 & 2.964 & 2.485 & -1.926 & (0.076) & 50.8 \\
2020 & 362 & 2.954 & 2.475 & -1.923 & (0.075) & 48.0 \\
\bottomrule\end{tabular}
\begin{tablenotes}\footnotesize
\item \emph{Notes:} Each row describes the occupations with a published mean gross hourly pay in that ASHE release on the SOC 2010 basis, of the 366 occupations in the panel. Log pay is the log of nominal mean gross hourly pay. The median automation probability across the 366 occupations is 0.468. The gradient is the slope from a cross-sectional regression of log pay on the automation probability in that release, with heteroskedasticity-consistent standard errors in parentheses. The final column is the share of employee jobs in above-median-risk occupations among the 309 occupations with a published jobs count in all seven releases.
\item \emph{Source:} ASHE Table 14, 2014--2020 releases, revised editions; ONS (2019) probability of automation.
\end{tablenotes}\end{threeparttable}\end{table}